\section{Related work}\label{sec:related}
\label{sec:related_work}

\subsection{Reinforcement learning for bus holding control}
\label{subsec:rl_holding}

The application of reinforcement learning to real-time bus holding control has attracted considerable attention in recent years, driven by the promise of adaptive policies that can respond to stochastic passenger demand and traffic conditions. \citet{alesiani2018reinforcement} formulated holding for high-frequency services as a value-based RL problem, showing that a Q-learning controller could reduce headway dispersion compared to schedule-based rules. \citet{wang2020real} subsequently studied multi-agent deep RL for the bus-bunching problem, casting each bus as an independent agent with a shared policy network and demonstrating that learned holding actions stabilize headways across an entire corridor without explicit communication. The shared-parameter paradigm common to these works scales gracefully with fleet size and implicitly encourages cooperative behavior through shared gradient updates.

Despite these advances, existing RL-based holding controllers share a fundamental limitation: they operate under a fixed, externally supplied timetable. The planned departure times and target headways are treated as immutable inputs, meaning that the holding policy has no mechanism to communicate execution-level feedback---such as systematic schedule infeasibility or recurrent bunching patterns---back to the planning stage. When the timetable itself is poorly calibrated to actual operating conditions, even an optimal holding policy can only mitigate, not resolve, the resulting service irregularity. TransitDuet addresses this gap by embedding the holding controller as the low-level agent within a bi-level framework, where real-time execution performance directly shapes the timetable through Temporal Advantage Propagation.

\subsection{Bus timetable optimization}
\label{subsec:timetable_opt}

Timetable design has traditionally been approached through mathematical programming and metaheuristic search. \citet{ibarra2015planning} formulated the problem as a mixed-integer program that jointly determines departure times and vehicle assignments to minimize a weighted combination of passenger waiting time and operator cost, solving it via branch-and-bound with valid inequalities. \citet{gkiotsalitis2021bus} proposed a bi-level model in which the upper level selects headways to minimize total passenger cost while the lower level solves a vehicle scheduling problem, solved using an adaptive large neighborhood search. Frequency-setting variants, such as the work of \citet{cats2019frequency}, optimize the number of departures per period rather than individual departure times, and typically rely on assignment models to estimate passenger flows under each candidate frequency vector.

A common thread in this literature is the reliance on static demand forecasts and deterministic travel-time assumptions. Timetables are optimized offline against expected conditions and then handed off to operations without a feedback loop. When actual conditions deviate---due to traffic incidents, demand surges, or driver variability---the timetable may become infeasible, and holding controllers are left to absorb the mismatch. Moreover, the combinatorial structure of these formulations limits scalability: solution times grow rapidly with the number of trips and stops, making re-optimization during the operating day impractical. TransitDuet replaces this open-loop paradigm with a closed-loop one. The high-level agent learns a timetable policy that maps aggregate state features to departure-time adjustments, and it receives gradient signals from low-level execution via the cross-level coupling channels described in Section~\ref{sec:tap}, enabling the planner to adapt its decisions based on how well the resulting schedule can actually be executed.

\subsection{Hierarchical and multi-timescale reinforcement learning}
\label{subsec:hierarchical_rl}

Hierarchical reinforcement learning decomposes complex tasks into temporally extended sub-problems. The options framework of \citet{sutton1999between} introduced the notion of temporally abstract actions, each consisting of an initiation set, an intra-option policy, and a termination condition, enabling planning over multiple timescales within a single MDP. Feudal networks \citep{vezhnevets2017feudal} operationalize this idea with a manager that sets sub-goals in a learned latent space and a worker that executes primitive actions to reach them, training both levels end-to-end with the manager's policy gradient modulated by the worker's goal-reaching success. HIRO \citep{nachum2018data} further improved sample efficiency by introducing off-policy corrections that relabel high-level goals in hindsight, making hierarchical training practical in continuous-control domains. Beyond goal-conditioned hierarchies, \citet{lee2020reinforcement} and \citet{metelli2020control} studied the multi-timescale setting more broadly, analyzing the interaction between agents operating at different decision frequencies and establishing regret bounds for the resulting non-stationary sub-problems.

While these methods provide powerful abstractions for temporal hierarchy, most assume either synchronous execution---where the high-level decision triggers immediately before a fixed-length sequence of low-level steps---or a shared state space across levels. In transit operations, the planning and control timescales are fundamentally asynchronous: the timetable agent acts once per planning horizon (e.g., every 15--30 minutes), while the holding agent acts at every stop visit (every 1--3 minutes), and the two levels observe different state representations at non-aligned decision epochs. TransitDuet handles this asynchrony through three coupling channels (Section~\ref{sec:tap}): a state channel exposing lower-level holding statistics to the upper policy, an action channel through which the upper offset directly perturbs the lower's environment, and a reward channel that backfills each dispatch's reward with a hindsight credit derived from the realised inter-dispatch gap. We initially attempted a literal asynchronous extension of RoboDuet-style cross-advantage augmentation, but the per-step advantage scales of the upper and lower policies differed substantially across a trip lifecycle and the resulting training was unstable; the reward-channel design described in this paper proved markedly more reliable.

\subsection{Constrained reinforcement learning}
\label{subsec:constrained_rl}

Many practical RL deployments must satisfy constraints beyond reward maximization. Lagrangian relaxation is the most widely used approach: \citet{tessler2019reward} proposed reward-constrained policy optimization, which augments the policy gradient with dual variables corresponding to cost constraints and performs primal-dual updates to steer the policy toward feasibility. \citet{stooke2020responsive} improved upon this with PID-based Lagrange multiplier updates that reduce oscillation and speed convergence in the constrained setting. These methods are effective for soft constraints, where occasional small violations are tolerable, but they offer no worst-case guarantees---the Lagrange multiplier may fail to grow fast enough to prevent persistent violations during training.

For settings that demand hard constraint satisfaction, \citet{miryoosefi2019reinforcement} introduced ApproPO, which formulates the constrained problem as an online convex optimization over the space of occupancy measures and applies a measurement-based projection oracle to enforce feasibility at every iterate. The $\theta$-OGD variant provides a practical implementation that projects policy parameters onto the feasible set using online gradient descent with approximate projections. In the multi-agent domain, \citet{pan2024roboduet} proposed RoboDuet for humanoid locomotion, where two coupled agents---an upper-body controller and a lower-body locomotor---exchange cross-advantage terms that capture how each agent's action quality is perceived under the other's value function, enabling tight coordination without centralized control.

TransitDuet combines elements from both lines of work to address the constraint structure of transit operations. Fleet-size feasibility is handled as a soft constraint via the $\theta$-OGD measurement-projection scheme of \citet{miryoosefi2019reinforcement}, which adaptively reweights an overshoot penalty in the upper-level reward; in expectation this drives the policy toward feasible fleet utilisation, but per-episode fluctuations under elastic fleet sampling and demand noise can still produce small, bounded overshoot. Headway uniformity is similarly handled as a soft constraint via a Lagrangian penalty with adaptive multiplier updates in the low-level holding policy. The cross-level coordination itself is managed by the trip-lifecycle reward channel and state-channel coupling described in Section~\ref{sec:tap}, which adapt the spirit of RoboDuet's synchronous cross-advantage to the asynchronous, multi-timescale regime that characterises the planning--control interface.
